% !TEX root = ../main.tex
\chapter{Reproduïbilitat i comandes}\label{app:reproducibility}
\fancyhead[RE]{\slshape \textsf{Apèndix \thechapter.\ Reproduïbilitat}}
\fancyhead[LO]{\slshape \textsf{\nouppercase \rightmark}}

Aquest apèndix descriu una ruta de reproducció gradual. Primer s'identifica el
repositori, l'estructura mínima, les dependències i el format de dades. Després
s'executa una prova curta que compila el codi C i processa una imatge amb la
ruta mesurada \(\lambda/Q\). Les campanyes completes de CSMR i Raspberry es
presenten al final com a reproducció avançada perquè requereixen dependències,
checkpoints o hardware específics.

\section{Repositori públic}

El repositori públic del projecte és:

\begin{center}
\url{https://github.com/Crist407/Implementacio-compressor-basat-en-xarxes-neurals-per-a-us-en-satel-lit}
\end{center}

El repositori conté el codi, pesos, configuració, dataset RAW públic i fonts de
la memòria. Els checkpoints complets, bundles Raspberry i entorns locals no
formen part de Git perquè són artefactes generats. El fitxer
\code{PUBLIC\_ARTIFACTS.md} resumeix quins elements són públics i quins es
regeneren amb scripts.

\section{Estructura pública del repositori}

\begin{table}[H]
\centering
\small
\begin{tabularx}{\textwidth}{lX}
\toprule
Directori & Contingut \\
\midrule
\code{src/c/} & Implementació C del compressor, descompressor, capes i
generadors de Q-map. \\
\code{src/python/analysis/} & Orquestradors i scripts d'anàlisi, inclosa la
ruta mesurada \(\lambda/Q\). \\
\code{src/python/reference/} & Referència Python/TensorFlow de SORTENY. \\
\code{scripts/raspberry/} & Benchmarks autocontinguts per Raspberry. \\
\code{scripts/docs/} & Generació de figures i compilació compatible de la
memòria. \\
\code{config/} & Llindars, calibració mínima \code{lambda005} i configuracions
de presets. \\
\code{weights/} & Pesos exportats del model necessaris per a la ruta C. \\
\code{data/} & Dataset RAW de 120 retalls, mostra mínima i instruccions de
visualització. \\
\code{docs/informe\_final/} & Fonts \LaTeX, figures i bibliografia de la
memòria. \\
\bottomrule
\end{tabularx}
\caption{Estructura del repositori públic utilitzada per reproduir els exemples
d'aquest apèndix.}
\label{tab:public_repo_structure}
\end{table}

\section{Instal·lació mínima}

La prova bàsica requereix un compilador C, \code{make}, Python 3 i NumPy. La
ruta CSMR necessita TensorFlow i TensorFlow Compression, però no són necessaris
per a la prova inicial.

\begin{verbatim}
python3 -m venv .venv
source .venv/bin/activate
pip install -r requirements.txt
\end{verbatim}

El repositori inclou una mostra RAW, el dataset de 120 retalls i la calibració
necessària per executar la ruta pública:

\begin{verbatim}
data/T31TCG_20230907T104629_5.8_512_512_2_1_0.raw
data/Sentinel2A_crop_test/
config/fq_calibration_lambda005.tsv
\end{verbatim}

\section{Dataset i visualització RAW}

Cada RAW té $8\times512\times512$ valors \code{uint16}, és a dir,
4.194.304 bytes. Els fitxers no tenen capçalera i estan guardats en ordre BSQ.
Per obrir-los amb Fiji/ImageJ cal utilitzar \code{File -> Import -> Raw...} i
configurar:

\begin{itemize}[leftmargin=*]
  \item tipus d'imatge: \code{16-bit Unsigned};
  \item amplada i alçada: $512\times512$;
  \item nombre d'imatges: 8;
  \item offset i separació entre imatges: 0;
  \item ordre de bytes: little-endian.
\end{itemize}

Per inspecció visual en color natural aproximat, les bandes recomanades són B04
com a vermell, B03 com a verd i B02 com a blau. El fitxer \code{data/README.md}
resumeix el format i els passos de visualització.

També es pot generar un paquet de quicklook sense executar el còdec:

\begin{verbatim}
python3 src/python/utils/manual_roi_quicklook.py \
  --input data/T31TCG_20230907T104629_5.8_512_512_2_1_0.raw \
  --output-dir output/checkpoints/quicklook_example \
  --pattern empty \
  --skip-qmap
\end{verbatim}

\section{Compilació C}

La compilació de la ruta C principal és:

\begin{verbatim}
make MODE=release OMP=1
make MODE=release OMP=1 test_ops
\end{verbatim}

Els binaris utilitzats per les proves d'aquest apèndix són:
\code{sorteny\_compressor}, \code{sorteny\_decompressor},
\code{sorteny\_fq\_qmap} i \code{sorteny\_semantic\_qmap}. El test
\code{test\_ops} comprova operadors interns del decoder abans d'executar una
campanya completa.

\section{Reproducció ràpida}

Després d'instal·lar dependències i compilar el C, la prova mínima processa una
imatge amb el mode \code{focus\_bgq128\_measured}:

\begin{verbatim}
python3 src/python/analysis/run_lambda005_measured_quality_route.py \
  --output-dir output/checkpoints/quickstart_measured_route \
  --modes focus_bgq128_measured \
  --threads 4
\end{verbatim}

Una execució correcta crea el directori
\code{output/checkpoints/quickstart\_measured\_route}. Aquesta prova valida el
camí essencial: entrada RAW, còdec C, error base, càlcul de qualitat local,
Q-map i reconstrucció final.

\section{Comprovació de sortides}

La prova ràpida ha de generar:

\begin{itemize}[leftmargin=*]
  \item Q-maps de 1.024 bytes;
  \item bitstreams C llegibles pel descompressor;
  \item reconstruccions RAW de 4.194.304 bytes;
  \item \code{metrics.csv}, \code{metrics.json} i \code{run\_meta.json}.
\end{itemize}

Les invariants principals es comproven amb:

\begin{verbatim}
CHECKPOINT=output/checkpoints/quickstart_measured_route

find "$CHECKPOINT/qmaps" -name '*.raw' -printf '%p %s\n'
find "$CHECKPOINT/reconstructions" -name '*.raw' -printf '%p %s\n'
sed -n '1,20p' "$CHECKPOINT/metrics.csv"
\end{verbatim}

\section{Comprovació dels extractes públics}

Les xifres principals de la memòria també es poden inspeccionar sense regenerar
els checkpoints complets. El directori \code{docs/informe\_final/evidence/}
conté extractes CSV petits amb els resums que alimenten el registre
d'evidències de l'Apèndix~\ref{app:evidence}.

\begin{verbatim}
ls docs/informe_final/evidence
sed -n '1,10p' docs/informe_final/evidence/csmr_policy_summary.csv
sed -n '1,10p' docs/informe_final/evidence/q204_same_mask_summary.csv
sed -n '1,10p' docs/informe_final/evidence/raspberry_operational_costs.csv
\end{verbatim}

El càlcul següent comprova tres valors destacats del cos principal: l'estalvi de
\code{focus\_bgq128} contra \code{q204}, el cost de la ruta amb error mesurat i
el cost relatiu de generar un Q-map semàntic.

\begin{verbatim}
python3 - <<'PY'
import csv

with open("docs/informe_final/evidence/csmr_policy_summary.csv",
          newline="") as f:
    rows = {r["policy"]: r for r in csv.DictReader(f)}
q204 = float(rows["q204"]["tfci_bps_mean"])
focus = float(rows["focus_bgq128"]["tfci_bps_mean"])
print("estalvi_focus_vs_q204_pct =", 100 * (1 - focus / q204))

with open("docs/informe_final/evidence/raspberry_operational_costs.csv",
          newline="") as f:
    costs = {r["metric"]: float(r["value"]) for r in csv.DictReader(f)}
print("ruta_error_mesurat_s =", costs["measured_error_route"])

with open("docs/informe_final/evidence/qmap_cost_by_type.csv",
          newline="") as f:
    qmap = {r["command_type"]: r for r in csv.DictReader(f)}
print("qmap_semantic_pct_compress =",
      qmap["semantic"]["qmap_pct_of_compress_mean"])
PY
\end{verbatim}

Aquesta comprovació valida les xifres agregades publicades. La regeneració de
tota una campanya continua requerint les rutes avançades de CSMR o Raspberry
descrites més avall.

\section{Ruta mesurada \texorpdfstring{$\lambda/Q$}{lambda/Q}}

La ruta mesurada implementa la cadena:

\begin{center}
\fbox{\begin{minipage}{0.92\textwidth}
\centering
compressió/descompressió base $\rightarrow$ error base $\rightarrow$
fórmules fixed-quality $\rightarrow$ $\lambda/Q$ $\rightarrow$ Q-map
$\rightarrow$ compressió final $\rightarrow$ mètriques
\end{minipage}}
\end{center}

Per executar una prova amb dos modes:

\begin{verbatim}
OUT=output/checkpoints/lambda005_measured_quality_route_smoke

python3 src/python/analysis/run_lambda005_measured_quality_route.py \
  --output-dir "$OUT" \
  --modes global_target_measured focus_bgq128_measured \
  --threads 4
\end{verbatim}

Per executar tots els modes implementats per defecte:

\begin{verbatim}
OUT=output/checkpoints/lambda005_measured_quality_route_full

python3 src/python/analysis/run_lambda005_measured_quality_route.py \
  --output-dir "$OUT" --threads 4
\end{verbatim}

Els modes suportats són \code{global\_target\_measured},
\code{adaptive\_s8\_measured}, \code{focus\_bgq128\_measured},
\code{target\_focus\_bgq128\_measured},
\code{preserve\_roi\_q240\_bgq128} i
\code{preserve\_roi\_q255\_bgq128}. Els detalls de configuració de cada mode es
descriuen a l'Apèndix~\ref{app:config}.

\section{Ruta precalibrada}

La ruta precalibrada utilitza informació obtinguda fora de la placa per generar
Q-maps sense repetir la passada base de compressor/descompressor. És útil per
proves ràpides, comparatives i escenaris en què la calibració sigui
representativa.

Q-map uniforme:

\begin{verbatim}
CAL=config/fq_calibration_lambda005.tsv

./sorteny_fq_qmap \
  --calibration "$CAL" \
  --target-from-q 204 \
  --output-qmap output/checkpoints/q204.raw
\end{verbatim}

Q-map semàntic amb preset i política:

\begin{verbatim}
CAL=config/fq_calibration_lambda005.tsv
IN=data/T31TCG_20230907T104629_5.8_512_512_2_1_0.raw

./sorteny_semantic_qmap \
  --input "$IN" \
  --bands 8 --height 512 --width 512 \
  --band-layout sentinel2-8 --preset cloud_avoid --threshold 0.50 \
  --semantic-policy focus --foreground-boost 16 --background-q 128 \
  --calibration "$CAL" \
  --output-qmap output/checkpoints/cloud_avoid_focus_bgq128.raw \
  --summary-tsv output/checkpoints/cloud_avoid_focus_bgq128.tsv
\end{verbatim}

\section{Reproducció avançada: CSMR i bitrate real}

El bitrate real de SORTENY es mesura amb la ruta Python/TensorFlow Compression
del model. Els auditors CSMR converteixen cada Q-map en un array de
$\lambda=Q\cdot0,05/255$ i executen la referència amb \code{quality\_array}.
Aquestes campanyes requereixen dependències Python addicionals i els
checkpoints de selecció corresponents.

\begin{verbatim}
SEL=output/checkpoints/20260620_lambda005_full_csmr_case_selection
OUT=output/checkpoints/20260623_lambda005_csmr_selected_cases_bitrate_final

python3 src/python/analysis/audit_csmr_selected_cases_from_full.py \
  --selection-checkpoint "$SEL" \
  --lambda-max 0.05 --output-dir "$OUT"
\end{verbatim}

\begin{verbatim}
THR=output/checkpoints/20260625_lambda005_threshold_optimization
OUT=output/checkpoints/20260626_lambda005_csmr_experimental_threshold_modes

python3 src/python/analysis/audit_csmr_experimental_threshold_modes.py \
  --threshold-checkpoint "$THR" \
  --lambda-max 0.05 --output-dir "$OUT"
\end{verbatim}

\begin{verbatim}
OUT=output/checkpoints/20260627_lambda005_csmr_preserve_roi_policy

python3 src/python/analysis/audit_csmr_preserve_roi_policy.py \
  --lambda-max 0.05 --output-dir "$OUT"
\end{verbatim}

Aquestes proves són més costoses que la ruta C-only i depenen de l'entorn
Python/TensorFlow. Per això el cos de la memòria separa clarament proxies C de
bitrate real \code{.tfci}.

\section{Reproducció avançada: Raspberry}

Els benchmarks Raspberry requereixen una placa configurada i els bundles
autocontinguts generats per al projecte. La prova del còdec optimitzat és:

\begin{verbatim}
nohup env INCLUDE_EXPERIMENTAL=1 THREADS_LIST="4" \
  scripts/raspberry/run_lambda005_bundle_benchmark.sh \
  > rpi_optimized_lambda005_t4.log 2>&1 &
\end{verbatim}

El cost de generar Q-maps sense executar el còdec s'aïlla amb:

\begin{verbatim}
nohup env INCLUDE_EXPERIMENTAL=1 REPEATS=3 \
  scripts/raspberry/run_lambda005_qmap_cost_benchmark.sh \
  > rpi_qmap_cost_lambda005.log 2>&1 &
\end{verbatim}

La integritat mínima d'un checkpoint Raspberry es comprova amb:

\begin{verbatim}
find "$CHECKPOINT" -name DONE | wc -l
find "$CHECKPOINT" -name run_meta.tsv | wc -l
find "$CHECKPOINT" -name '*.time.txt' | wc -l
find "$CHECKPOINT" -name qmap.raw -size 1024c | wc -l
tar -tzf "${CHECKPOINT}.tar.gz" >/dev/null
sha256sum "${CHECKPOINT}.tar.gz"
\end{verbatim}

\section{Figures i compilació de la memòria}

Les figures derivades del cos principal es regeneren amb:

\begin{verbatim}
python3 scripts/docs/build_final_report_figures.py
\end{verbatim}

La memòria es compila amb el flux compatible, sense modificar la plantilla
institucional original:

\begin{verbatim}
scripts/docs/build_informe_final_compat.sh
\end{verbatim}

Una compilació correcta no ha de deixar cites, referències ni figures sense
resoldre.

\section{Verificacions locals finals}

Abans de considerar una reproducció local com a vàlida, s'han d'executar com a
mínim:

\begin{verbatim}
SCRIPT=src/python/analysis/run_lambda005_measured_quality_route.py
python3 -m py_compile "$SCRIPT"
make MODE=release OMP=1
make MODE=release OMP=1 test_ops
git diff --check
\end{verbatim}

Les campanyes completes generen més artefactes que els inclosos a Git. La seva
procedència, mida, hash i relació amb les taules i figures del cos principal es
documenten a l'Apèndix~\ref{app:evidence}.
